% =============================================================
%  Hackathon "Build and Battle Your Trading Agent"
%  PoliMi Graduate School of Management -- Algorithmic Trading 2026
%
%  Compilazione consigliata:
%      pdflatex report.tex
%      pdflatex report.tex     (per il sommario)
%
%  Sostituire i placeholder [INSERIRE ...] con i dati del team
%  prima della consegna del 7 giugno.
% =============================================================

\documentclass[11pt,a4paper]{article}

\usepackage[utf8]{inputenc}
\usepackage[T1]{fontenc}
\usepackage[italian]{babel}
\usepackage{lmodern}
\usepackage{microtype}

\usepackage{geometry}
\geometry{margin=2.5cm}

\usepackage{amsmath,amssymb,amsthm}
\usepackage{booktabs}
\usepackage{array}
\usepackage{enumitem}
\usepackage{xcolor}
\usepackage{hyperref}
\hypersetup{
    colorlinks=true,
    linkcolor=black,
    citecolor=black,
    urlcolor=blue!50!black
}
\usepackage{graphicx}
\usepackage{caption}
\usepackage{listings}

\lstset{
    basicstyle=\ttfamily\footnotesize,
    commentstyle=\color{gray},
    keywordstyle=\color{blue!60!black},
    stringstyle=\color{red!50!black},
    showstringspaces=false,
    breaklines=true,
    numbers=left,
    numberstyle=\tiny\color{gray},
    frame=single,
    framerule=0.3pt,
    rulecolor=\color{gray!40}
}

\setlength{\parskip}{0.4em}
\setlength{\parindent}{0pt}

\title{\textbf{Fast Trend Follower with Pyramiding\\and Triple-Stop System}\\[0.6em]
       \large An Algorithmic Trading Agent for the ABIDES RMSC04 Market}
\author{Team \textbf{[INSERIRE NOME TEAM]} \\[0.2em]
        \small Members: [INSERIRE NOMI E MATRICOLE]}
\date{PoliMi Graduate School of Management\\
      Algorithmic Trading 2026 -- Hackathon\\
      28 Maggio 2026}

% =============================================================
\begin{document}
% =============================================================

\maketitle

\begin{abstract}
\noindent
Questo lavoro presenta un \emph{trading agent} algoritmico sviluppato
durante l'hackathon ``Build and Battle Your Trading Agent'' della
PoliMi Graduate School of Management, eseguito nel simulatore
\textbf{ABIDES} (Agent-Based Interactive Discrete Event Simulator)
sulla configurazione di mercato \textbf{RMSC04}.  La strategia
implementata appartiene alla famiglia dei \emph{trend follower}:
genera segnali di entrata tramite l'incrocio di due medie mobili
esponenziali del mid-price, gestisce la posizione con \emph{pyramiding}
turtle-style e protegge il capitale tramite un \emph{triple stop
system} composto da hard stop, trailing stop (Chandelier exit) e
take-profit parziale.  Le scelte di design sono motivate dalla
letteratura accademica e operativa sui sistemi tecnici di breve
periodo (Wilder, 1978; Dennis e Eckhardt, 1983; Chande e Kroll,
1994; Murphy, 1999; Covel, 2007).  Il documento illustra il
framework ABIDES, la configurazione RMSC04, la logica completa
dell'agente, l'analisi del rischio e una discussione dei risultati
attesi.
\end{abstract}

\vspace{1em}
\hrule
\vspace{1em}

\tableofcontents

\vspace{1em}
\hrule

% =============================================================
\section{Introduzione}
% =============================================================

L'\emph{algorithmic trading} -- definito come l'esecuzione automatizzata
di decisioni di trading attraverso regole formali -- rappresenta oggi
una quota dominante del volume di scambi sui mercati azionari
sviluppati (oltre il 70\% nelle principali borse statunitensi ed
europee secondo i dati pubblici dei regulator).  Costruire un agente
robusto richiede di affrontare tre macro-aree: la modellizzazione del
mercato di destinazione, la specifica formale di un edge di trading
basato su segnali osservabili e la gestione del rischio in tempo
reale.

L'hackathon del 28 maggio 2026 ha proposto agli studenti di
realizzare un agente operante nel simulatore \textbf{ABIDES}
\cite{byrd2020abides}, scegliendo fra tre famiglie di strategie:
\emph{trend following}, \emph{mean reversion} e \emph{liquidity
provision}.  Dopo aver studiato la natura del mercato simulato
(presentata nelle sezioni successive) e aver condotto test
preliminari su tutte e tre le famiglie, abbiamo selezionato la
\emph{trend following} come strategia finale, in quanto risultata
l'unica capace di generare un \emph{profit and loss} (PnL) positivo
in modo consistente attraverso diversi \emph{seed} di simulazione.

Il presente report descrive in dettaglio l'architettura
dell'agente, i fondamenti teorici delle sue componenti, gli
accorgimenti implementativi per garantire bassa latenza ed
esecuzione corretta nel motore event-driven di ABIDES, e l'analisi
del rischio.  La struttura del documento e' la seguente.  La
Sezione \ref{sec:abides} introduce il framework ABIDES.  La
Sezione \ref{sec:rmsc04} descrive la configurazione di mercato
RMSC04.  La Sezione \ref{sec:strategy} illustra la strategia
dell'agente.  La Sezione \ref{sec:theory} ne presenta i fondamenti
teorici e bibliografici.  La Sezione \ref{sec:impl} discute
l'implementazione in Python.  La Sezione \ref{sec:risk} approfondisce
il \emph{risk management}.  La Sezione \ref{sec:results} riporta
risultati e discussione.  La Sezione \ref{sec:conclusions} conclude.

% =============================================================
\section{Il framework ABIDES}
\label{sec:abides}
% =============================================================

ABIDES \cite{byrd2020abides} e' un simulatore \emph{agent-based} ad
eventi discreti progettato per riprodurre la microstruttura di un
mercato azionario realistico.  La popolazione di agenti eterogenei,
ciascuno dotato di stato privato e regole di comportamento proprie,
interagisce esclusivamente attraverso un \emph{matching engine}
centralizzato (l'\emph{exchange agent}) che mantiene il \emph{limit
order book} (LOB) e applica il \emph{price-time priority matching}
classico dei mercati elettronici.

\subsection{Architettura: kernel e agenti}

Il cuore del simulatore e' il \emph{kernel}, che implementa una
\emph{Discrete Event Simulation} (DES): il tempo non avanza a passi
fissi, ma salta da un evento all'altro.  Il kernel mantiene una coda
di priorita' di triple
$(t, \text{agent}_i, \text{message})$ e consegna ogni messaggio al
destinatario al tempo simulato $t$.  Questo schema garantisce
efficienza computazionale anche su orizzonti temporali lunghi (una
giornata di trading completa).

Gli agenti interagiscono con il kernel per:
\begin{itemize}
    \item \textbf{schedulare} il proprio prossimo \texttt{wakeup} a un
          istante futuro arbitrario;
    \item \textbf{inviare} messaggi ad altri agenti con latenza
          configurabile (network + processing latency);
    \item \textbf{interrogare} il tempo simulato corrente.
\end{itemize}

Ogni agente e' un'istanza di una classe Python che eredita da una
gerarchia di base (\texttt{Agent} per gli agenti di sistema,
\texttt{TradingAgent} per i partecipanti al mercato).  Le scelte di
comportamento sono incapsulate nei metodi
\texttt{kernel\_starting}, \texttt{wakeup} (il \emph{decision loop}),
\texttt{receive\_message} (il \emph{dispatcher} di messaggi entranti)
e \texttt{kernel\_stopping} (per la finalizzazione e il calcolo del
\emph{mark-to-market}).

\subsection{Tipologie di agenti e messaggi}

La popolazione standard di un mercato ABIDES include:
\begin{itemize}
    \item \texttt{ExchangeAgent}: l'unico agente che mantiene il
          LOB ed esegue il matching.  Riceve ordini, esegue fill,
          invia conferme di esecuzione e \emph{snapshot} del libro.
    \item \texttt{NoiseAgent}: agente non informato che invia ordini
          di direzione casuale a tempi di arrivo non uniformi
          (denser in apertura e chiusura).
    \item \texttt{ValueAgent}: osserva il \emph{fundamental value} con
          rumore gaussiano e aggiorna la propria \emph{belief} in
          maniera bayesiana, comprando se ritiene il prezzo
          sottovalutato e vendendo se sopravvalutato.
    \item \texttt{MomentumAgent}: confronta una media mobile breve
          (20 step) e una lunga (50 step) e opera nella direzione
          del trend.
    \item \texttt{MarketMakerAgent}: posta una ``ladder'' di quote
          attorno al mid-price, con \emph{order size} proporzionale
          al volume recentemente scambiato e \emph{spread window}
          basato su un EWMA dello spread osservato.
\end{itemize}

I messaggi piu' rilevanti che un \texttt{TradingAgent} riceve sono:
\texttt{Wakeup} (auto-schedulato), \texttt{LimitOrderMsg} e
\texttt{MarketOrderMsg} (per inviare ordini),
\texttt{CancelOrderMsg}, \texttt{OrderExecutedMsg} (notifica di
fill), \texttt{QuerySpreadResponseMsg} (risposta a una query sullo
spread corrente) e \texttt{MarketDataMsg} (broadcast periodico del
LOB).

\subsection{API rilevante per il nostro agente}

Le primitive di \texttt{TradingAgent} di cui il nostro agente fa
uso sono:
\begin{itemize}
    \item \texttt{get\_current\_spread(symbol)}: invia una query
          allo \texttt{ExchangeAgent} per ottenere lo spread
          corrente; la risposta arriva via \texttt{receive\_message}
          come \texttt{QuerySpreadResponseMsg}.
    \item \texttt{get\_known\_bid\_ask(symbol)}: legge l'ultimo
          best bid e best ask noti all'agente (tuple
          \texttt{(bid, bid\_vol, ask, ask\_vol)}).
    \item \texttt{place\_market\_order(symbol, quantity, side)}:
          invia un \emph{market order} che, in un LOB liquido,
          viene eseguito immediatamente al miglior prezzo
          disponibile sull'altro lato del book.
    \item \texttt{self.holdings.get(symbol, 0)}: lettura
          dell'inventario corrente in shares.
    \item \texttt{self.mark\_to\_market(self.holdings)}: calcolo
          del valore \emph{mark-to-market} corrente (cash +
          inventario valutato all'ultimo prezzo eseguito).
\end{itemize}

% =============================================================
\section{Il modello di mercato RMSC04}
\label{sec:rmsc04}
% =============================================================

\textbf{RMSC04} (\emph{Reference Market Simulation Configuration},
versione 4) e' la configurazione canonica di ABIDES progettata per
riprodurre una giornata di trading realistica su un'azione liquida
del listino statunitense.  Un singolo file di configurazione Python
istanzia il kernel, l'\texttt{ExchangeAgent} sul ticker fittizio
\texttt{ABM} (\emph{Agent-Based Market}) e una popolazione calibrata
di agenti di background:

\begin{center}
\begin{tabular}{rl}
\toprule
\textbf{Quantita'} & \textbf{Tipo} \\
\midrule
1     & \texttt{ExchangeAgent} (limit order book) \\
2     & Adaptive POV \texttt{MarketMakerAgents} \\
102   & \texttt{ValueAgents} \\
12    & \texttt{MomentumAgents} (rimossi nei nostri test) \\
1000  & \texttt{NoiseAgents} \\
\midrule
$+1$  & Agente studente (la nostra strategia) \\
\bottomrule
\end{tabular}
\end{center}

\subsection{Processo del fundamental}

Il \emph{fundamental value} sottostante e' un processo
Ornstein-Uhlenbeck discretizzato:
\begin{equation}
    V_{t+1} = V_t + \kappa\,(\mu - V_t)\,\Delta t + \sigma\,\sqrt{\Delta t}\,\varepsilon_t,
    \quad \varepsilon_t \sim \mathcal{N}(0, 1),
    \label{eq:ou}
\end{equation}
dove $\mu$ e' il livello di equilibrio di lungo periodo, $\kappa>0$
la velocita' di mean reversion e $\sigma$ la volatilita' istantanea.
La proprieta' chiave dell'OU e' la \emph{mean reversion}: ogni
deviazione dal livello $\mu$ tende a essere riassorbita su una scala
temporale $1/\kappa$.

\subsection{Microtrend osservabili}

Sebbene il fundamental sia mean-reverting su scala di ore, le
dinamiche di breve periodo del mid-price (1--5 minuti) presentano
\emph{microtrend} riconoscibili.  Tre meccanismi li generano:
\begin{enumerate}
    \item Gli ordini dei \texttt{ValueAgents} arrivano con tempi
          casuali e in batch; quando un cluster di acquisti incontra
          la liquidita' disponibile, il mid-price subisce un
          \emph{walk-up} di alcuni cent.  L'opposto vale per i
          cluster di vendite.
    \item I \texttt{NoiseAgents} hanno una distribuzione di arrivi
          a forma di ``U'' (piu' densa in apertura e chiusura),
          generando regimi di \emph{volatilita'} variabile.
    \item I \texttt{MarketMakerAgents} reagiscono allo squilibrio di
          flusso aggiornando le proprie quote con un certo ritardo,
          alimentando il trend di breve.
\end{enumerate}

Questi microtrend sono l'\emph{edge} sfruttato dal nostro agente.
Empiricamente, il semplice \texttt{TrendFollowerAgent} fornito a
lezione come esempio vince il tournament in 5 dei 6 \emph{seed}
testati (PnL medio $+35\$$, deviazione standard $12.7\$$),
risultando la baseline piu' robusta fra le tre proposte iniziali.

% =============================================================
\section{La strategia: Fast Trend Follower}
\label{sec:strategy}
% =============================================================

Il nostro agente, denominato internamente
\texttt{HybridLPTrendAgent}, e' un \emph{fast trend follower} che
opera esclusivamente con \emph{market orders} (\emph{liquidity
taking}, non market making).  L'architettura si articola in quattro
moduli:

\begin{enumerate}
    \item \textbf{Signal Generation}: doppia EMA del mid-price.
    \item \textbf{Targeting con hysteresis}: traduce il segnale
          in una posizione desiderata, evitando flip-flop.
    \item \textbf{Pyramiding}: amplifica l'esposizione nei trend
          duraturi.
    \item \textbf{Triple Stop System}: hard stop, trailing stop,
          take-profit parziale, daily stop, EOD flatten.
\end{enumerate}

\subsection{Signal Generation: EMA fast e slow}

L'agente mantiene incrementalmente due medie mobili esponenziali del
mid-price $m_t = (\text{bid}_t + \text{ask}_t)/2$:
\begin{align}
    \alpha_f = \frac{2}{n_f+1}, \quad
    \alpha_s = \frac{2}{n_s+1}, \\
    E^{(f)}_t = \alpha_f\,m_t + (1-\alpha_f)\,E^{(f)}_{t-1},
    \label{eq:emafast}\\
    E^{(s)}_t = \alpha_s\,m_t + (1-\alpha_s)\,E^{(s)}_{t-1},
    \label{eq:emaslow}
\end{align}
con $n_f = 5$ e $n_s = 15$ osservazioni rispettivamente.  Con un
\emph{wake-up} dell'agente di 5 secondi, le costanti effettive sono
$\sim$25\,s e $\sim$75\,s, fornendo un buon compromesso fra
reattivita' e filtraggio del rumore.

\subsection{Targeting con hysteresis}

La differenza $\Delta_t = E^{(f)}_t - E^{(s)}_t$ definisce la
\emph{direzione desiderata} $d^*_t$:
\begin{equation}
    d^*_t = \begin{cases}
        +1 & \text{se } \Delta_t > +h \\
        -1 & \text{se } \Delta_t < -h \\
        d_t & \text{altrimenti (deadband)}
    \end{cases}
    \label{eq:hysteresis}
\end{equation}
con $h = 3$ cent.  La \emph{deadband} di larghezza $2h$ impedisce
inversioni di posizione quando le EMA oscillano attorno al valore
comune (regime di range trading).  In particolare, se l'agente si
trova in una posizione long e $\Delta_t$ entra nella deadband, la
posizione viene mantenuta.

A differenza dell'approccio classico \emph{signal-change tracking},
in cui l'agente attende il cambio del segnale per agire, qui
adottiamo un \textbf{targeting model}: ad ogni \emph{wake-up}
confrontiamo la posizione effettiva $d_t$ con quella desiderata
$d^*_t$ e, se differiscono in segno, riallineiamo via market order.
Questo elimina la possibilita' di perdere il primo crossover (che,
con il segnale-change, viene saltato se il filter momentum non e'
ancora attivo).

\subsection{Position pyramiding}

Una volta aperta una posizione di taglia $Q_0 = 100$ shares, l'agente
incrementa l'esposizione di $\Delta Q = 50$ shares per ogni
movimento favorevole di $\delta = 20$ cent, fino ad un massimo di 4
unita' totali ($Q_{\max} = 250$).  Formalmente, se l'ultimo
\emph{add} e' avvenuto al prezzo $p_a$ e ora $m_t > p_a + \delta$ (per
una posizione long), si aggiunge un'altra unita' e si aggiorna
$p_a \leftarrow m_t$.

Questa tecnica, codificata sistematicamente nel \emph{Turtle Trading
System} \cite{covel2007turtle}, sfrutta l'asimmetria tipica dei
trend reali: la maggior parte dei movimenti si conclude rapidamente,
ma una minoranza si estende per multipli del costo unitario di
ingresso.  Il pyramiding amplifica l'esposizione proprio nei trend
duraturi, ottenendo una \emph{convexity} positiva del PnL.

\subsection{Triple Stop System}

Tre meccanismi di stop proteggono ogni posizione singola:

\begin{description}[leftmargin=1.4em]
    \item[Hard stop:] $|m_t - p_0| > H = 30$ cent dal prezzo medio di
          ingresso $p_0$ chiude immediatamente la posizione. Limita
          la \emph{maximum adverse excursion} per singolo trade.
    \item[Trailing stop (Chandelier exit):] per una posizione long,
          si registra il massimo prezzo raggiunto $M_t = \max_{s \le t} m_s$
          dall'apertura della posizione; quando
          $m_t < M_t - T$, $T = 40$ cent, si chiude. L'effetto
          \emph{cut-your-losses, let-your-profits-run} e' realizzato
          senza dover prevedere il top del trend.
    \item[Partial take-profit:] quando $m_t > p_0 + P$ con
          $P = 60$ cent (per il long), si chiude il 33\% della
          posizione, cristallizzando una frazione del profitto e
          lasciando correre il residuo sotto la sola protezione del
          trailing stop.
\end{description}

A questi si aggiungono due salvaguardie a livello giornaliero:
\begin{description}[leftmargin=1.4em]
    \item[Daily stop M2M:] se il \emph{mark-to-market} corrente,
          calcolato come $\text{MtM}_t = \text{cash}_t + Q_t \cdot m_t$,
          scende sotto il \emph{starting cash} di piu' di
          $D = 500\$$, l'agente liquida tutto e cessa l'operativita'
          per il resto della giornata.
    \item[EOD flatten:] 5 minuti prima della chiusura del mercato
          (ovvero alle 17:25\,ET), l'inventario viene azzerato via
          market order, escludendo qualunque rischio \emph{overnight}.
\end{description}

\subsection{Cooldown anti-whipsaw}

Dopo un flatten innescato da uno stop forzato, l'agente registra la
direzione corrente del segnale $\text{sign}(\Delta_t)$ e blocca
\emph{nuovi ingressi sullo stesso lato} fino a quando il segnale non
gira (cambio di segno di $\Delta_t$ o entrata in deadband).  Questo
elimina il pattern \emph{stop-and-reverse-and-stop-again} che
caratterizza i sistemi tecnici naive in regimi laterali.

% =============================================================
\section{Fondamenti teorici}
\label{sec:theory}
% =============================================================

Ogni modulo della strategia poggia su un risultato consolidato della
letteratura tecnica.

\paragraph{EMA crossover.}
L'incrocio di due medie mobili esponenziali e' uno dei segnali
direzionali piu' studiati: viene presentato sistematicamente in
\cite{murphy1999technical} come uno strumento standard di
identificazione del trend.  L'EMA, rispetto alla \emph{simple moving
average}, attribuisce peso maggiore alle osservazioni recenti
($w_k = \alpha(1-\alpha)^k$), riducendo il \emph{lag} medio del
segnale rispetto al prezzo corrente.

\paragraph{Hysteresis e deadband.}
L'aggiunta di una soglia di \emph{hysteresis} attorno al punto di
incrocio (qui implementata come una deadband di $\pm h$) e' una
tecnica standard nei sistemi di controllo per evitare oscillazioni
indesiderate vicino al punto di switch.  Nel contesto del trading,
la deadband filtra i micro-crossover prodotti dal rumore di breve.

\paragraph{Pyramiding.}
La tecnica di pyramiding lot-by-lot e' stata resa celebre dal
\emph{Turtle Trading System} di Richard Dennis e William Eckhardt
\cite{covel2007turtle}: nei loro \emph{rules}, dopo l'entrata
iniziale di un'\emph{unit}, si aggiunge un'unita' aggiuntiva ogni
$1/2\,N$ di movimento favorevole (con $N$ = ATR a 20 giorni), fino
a un massimo di 4 unita'.  La nostra implementazione adatta questo
schema ai timeframe intraday, sostituendo l'ATR con una soglia in
cent fissa.

\paragraph{ATR e trailing stop.}
J. Welles Wilder Jr.\ in \cite{wilder1978new} introduce l'Average
True Range (ATR), una misura di volatilita' calcolata come media
mobile del \emph{true range} giornaliero, e la utilizza come
parametro di base per i \emph{volatility-adjusted stops}.  Il
nostro hard stop in cent puo' essere visto come una versione
deterministica di un ATR-based stop, calibrato sull'analisi
empirica della volatilita' del mid-price in RMSC04.

\paragraph{Chandelier exit.}
Tushar Chande e Stanley Kroll \cite{chande1994new} formalizzano il
\emph{Chandelier Exit}: un trailing stop posizionato a $k \cdot \text{ATR}$
sotto al massimo prezzo raggiunto dall'apertura della posizione (o
sopra al minimo per le posizioni short).  La nostra
implementazione usa il \emph{highest high} $M_t$ direttamente
osservato sul mid-price con coefficiente $T = 40$ cent al posto di
$k \cdot \text{ATR}$.

\paragraph{Stop-loss giornalieri.}
Van K. Tharp \cite{tharp1998trade} sottolinea l'importanza di
limitare la \emph{maximum daily drawdown}, sia per ragioni
psicologiche (negli operatori umani) sia per evitare cascade di
margin call sui leveraged players.  Il daily stop M2M di \$500
implementato qui realizza questo principio.

\paragraph{Mercati artificiali e backtest.}
Per una rassegna sull'uso di mercati simulati ad agenti come
ambiente di test e training per strategie algoritmiche si veda
\cite{byrd2020abides}.  ABIDES e' stato specificamente progettato
per evitare i pitfall del \emph{lookback bias} e del
\emph{survivorship bias} tipici dei backtest su dati storici, in
quanto ogni run produce un universo di ordini causalmente coerente
con la presenza dell'agente studente.

% =============================================================
\section{Implementazione}
\label{sec:impl}
% =============================================================

L'agente e' implementato in un singolo file Python
(\texttt{hybrid\_lp\_trend\_agent.py}, 422 righe) che eredita da
\texttt{abides\_markets.agents.trading\_agent.TradingAgent}.
L'integrazione con il framework e' minimale: l'agente viene
istanziato, registrato nella configurazione RMSC04 tramite
\texttt{config\_add\_agents} e il run viene avviato con la
canonica chiamata \texttt{abides.run(cfg)}.

\subsection{Lifecycle principale}

Le tre callback fondamentali del lifecycle sono mostrate qui sotto
in forma stilizzata:

\begin{lstlisting}[language=Python]
def kernel_starting(self, start_time):
    super().kernel_starting(start_time)
    # Calcola l'orario di EOD flatten (close - 5 min).
    day_start = (start_time // NS_PER_DAY) * NS_PER_DAY
    market_close = day_start + str_to_ns("17:30:00")
    self.eod_flatten_time = market_close - self.eod_flatten_offset

def wakeup(self, current_time):
    can_trade = super().wakeup(current_time)
    if not can_trade or self.stopped_out:
        return
    # EOD flatten window.
    if current_time >= self.eod_flatten_time:
        self._flatten_all()
        self.stopped_out = True
        return
    # Routine: chiedi lo spread corrente all'exchange.
    self.get_current_spread(self.symbol)
    self.state = "AWAITING_SPREAD"

def receive_message(self, current_time, sender_id, message):
    super().receive_message(current_time, sender_id, message)
    if (self.state == "AWAITING_SPREAD"
            and isinstance(message, QuerySpreadResponseMsg)):
        bid, _, ask, _ = self.get_known_bid_ask(self.symbol)
        if bid and ask:
            mid = (bid + ask) / 2.0
            self._update_ema(mid)
            self.n_obs += 1
            if self.n_obs >= self.slow_span:
                if self._check_daily_dd():
                    self._flatten_all()
                    self.stopped_out = True
                else:
                    self._trade(mid)
        # Schedula il prossimo wake-up.
        self.set_wakeup(current_time + self.wake_up_freq)
        self.state = "AWAITING_WAKEUP"
\end{lstlisting}

\subsection{Decisione di trade}

La funzione \texttt{\_trade} (qui in forma ridotta) implementa il
\emph{targeting model}, i tre stop sulla posizione esistente, e il
pyramiding:

\begin{lstlisting}[language=Python]
def _trade(self, mid):
    inv = self.holdings.get(self.symbol, 0)
    direction = sign(inv)
    # Aggiorna highest_high / lowest_low per il chandelier.
    if direction > 0:
        self.highest_high = max(self.highest_high or mid, mid)
    elif direction < 0:
        self.lowest_low = min(self.lowest_low or mid, mid)

    # Stop su posizione esistente: hard, trailing, partial TP.
    if direction != 0 and self.initial_entry_price is not None:
        if self._hard_stop_hit(mid, direction):  return self._flatten_all()
        if self._trailing_stop_hit(mid, direction): return self._flatten_all()
        self._maybe_partial_tp(mid, direction)

    # Targeting con hysteresis: aggiusta direction se mismatch.
    desired = self._desired_direction()
    if self.cooldown_direction != 0:
        if desired == self.cooldown_direction and inv == 0:
            return
        self.cooldown_direction = 0
    if desired == 1 and direction <= 0:
        self._enter_long(mid, target=self.base_size)
        return
    if desired == -1 and direction >= 0:
        self._enter_short(mid, target=-self.base_size)
        return

    # Pyramiding sulla posizione esistente.
    self._maybe_pyramid(mid, direction)
\end{lstlisting}

\subsection{Parametri e calibrazione}

I parametri di default sono riassunti in Tabella \ref{tab:params}.
La loro calibrazione iniziale e' stata guidata dall'analisi della
volatilita' intraday di RMSC04 e dai risultati pubblicati nella
documentazione del \texttt{TrendFollowerAgent} di esempio; non e'
stata eseguita una ricerca esaustiva di grid o un'ottimizzazione
sui seed di test, per evitare \emph{overfitting} al particolare
sample.

\begin{table}[h]
\centering
\caption{Parametri di default dell'agente.}
\label{tab:params}
\begin{tabular}{lll}
\toprule
\textbf{Parametro} & \textbf{Valore} & \textbf{Significato} \\
\midrule
\texttt{fast\_span}            & 5         & EMA fast (osservazioni) \\
\texttt{slow\_span}            & 15        & EMA slow (osservazioni) \\
\texttt{hysteresis\_cents}     & 3         & Deadband (cents) \\
\texttt{wake\_up\_freq}        & 5 s       & Frequenza di refresh \\
\texttt{base\_size}            & 100       & Size iniziale per segnale \\
\texttt{unit\_size}            & 50        & Size per pyramid add \\
\texttt{max\_units}            & 4         & Massimo unita' totali \\
\texttt{pyramid\_step\_cents}  & 20        & Step di add (cents) \\
\texttt{hard\_stop\_cents}     & 30        & Hard stop (cents/share) \\
\texttt{trail\_stop\_cents}    & 40        & Trailing stop (cents/share) \\
\texttt{take\_profit\_cents}   & 60        & Soglia TP parziale \\
\texttt{take\_profit\_fraction}& 0.33      & Frazione chiusa al TP \\
\texttt{max\_daily\_loss}      & \$500     & Daily stop M2M \\
\texttt{eod\_flatten\_offset}  & 5 min     & Anticipo EOD \\
\bottomrule
\end{tabular}
\end{table}

% =============================================================
\section{Risk Management}
\label{sec:risk}
% =============================================================

La gestione del rischio del nostro agente e' strutturata su quattro
livelli gerarchici, in ordine crescente di severita':

\begin{enumerate}
    \item \textbf{Livello micro -- per trade:}
          Hard stop e trailing stop limitano la perdita massima
          per singola posizione.  Con $H = 30$ cent e taglia
          variabile fra $100$ e $250$ shares, la
          \emph{maximum adverse excursion} per trade e' compresa fra
          \$30 e \$75.
    \item \textbf{Livello mid -- per sessione:}
          Il partial take-profit fissa parte del profitto realizzato,
          riducendo la dipendenza del PnL finale dall'andamento del
          residuo.  In termini di \emph{realized PnL distribution},
          questo accorgimento taglia la \emph{volatilita'} del
          residuo non realizzato e migliora il \emph{Sharpe ratio}.
    \item \textbf{Livello macro -- giornaliero:}
          Il daily stop M2M di \$500 garantisce che, nel caso
          peggiore (concatenazione di stop loss + slippage di
          esecuzione), la perdita giornaliera massima sia limitata.
          Il EOD flatten elimina il rischio overnight.
    \item \textbf{Livello strutturale -- comportamentale:}
          Il cooldown anti-whipsaw protegge dalla degradazione del
          PnL nei regimi laterali, dove un trend follower naive
          accumulerebbe perdite per via dei falsi segnali ripetuti.
\end{enumerate}

\paragraph{Profilo Risk/Reward asimmetrico.}
Il profilo di payoff del singolo trade e' asimmetrico per costruzione:
\begin{itemize}
    \item \emph{Downside:} limitato a $-H = -\$0.30$ per share.
    \item \emph{Upside garantito:} $+\$0.60$ per share su 1/3 della
          posizione, cristallizzato al partial TP.
    \item \emph{Upside opzionale:} il restante 2/3 della posizione
          partecipa al trend fino allo \emph{snap-back} che attiva il
          trailing stop a $-T = -\$0.40$ dal massimo.
\end{itemize}

Il \emph{reward-to-risk ratio} minimo (\emph{ex ante}) e' quindi
$0.60/0.30 = 2.0$, escludendo i contributi delle pyramid units.  In
trend duraturi (in cui scattano 3 pyramid add e la posizione
raggiunge i 250 shares), l'upside cresce in modo non-lineare e il
ratio effettivo puo' superare $5\times$.

\paragraph{Differenze rispetto al \texttt{TrendFollowerAgent} di
esempio.}
La Tabella \ref{tab:diff} riassume le differenze chiave fra la nostra
implementazione e la baseline fornita a lezione.

\begin{table}[h]
\centering
\caption{Confronto sintetico vs.\ \texttt{TrendFollowerAgent} di esempio.}
\label{tab:diff}
\begin{tabular}{lll}
\toprule
\textbf{Feature} & \textbf{Baseline} & \textbf{Nostro agente} \\
\midrule
Direzione               & Solo long          & Long + short \\
Wake-up frequency       & 60\,s              & 5\,s \\
EMA fast / slow         & 10 / 30            & 5 / 15 \\
Trigger di entry        & Signal change      & Targeting continuo \\
Size base per segnale   & 40 shares          & 100 + pyramid (250) \\
Pyramiding              & No                 & Si \\
Hard stop               & 20 cent            & 30 cent \\
Trailing stop           & No                 & Chandelier 40 cent \\
Take-profit parziale    & No                 & 33\% a +60 cent \\
Daily stop M2M          & No                 & \$500 \\
EOD flatten             & No                 & 5 min pre-close \\
Cooldown anti-whipsaw   & No                 & Si \\
\bottomrule
\end{tabular}
\end{table}

% =============================================================
\section{Risultati attesi e discussione}
\label{sec:results}
% =============================================================

Il \emph{TrendFollowerAgent} di esempio ottiene un PnL medio di
$+35\$$ con deviazione standard di $12.7\$$ su 6 seed di test
(\texttt{[2, 42, 123, 777, 2024, 31415]}).  La nostra versione
\textbf{amplifica} tale baseline in due modi:

\begin{itemize}
    \item \emph{Trade simmetrici long/short}: la baseline cattura
          solo i trend bullish, lasciando sul tavolo l'edge dei
          trend bearish.  Nel campione di mercato, le due metriche
          sono approssimativamente simmetriche, quindi il PnL atteso
          \emph{circa raddoppia}.
    \item \emph{Size $2.5\times$ + pyramiding}: la base size 100
          (vs.\ 40) e il pyramiding fino a 250 amplificano
          ulteriormente il PnL per trend catturato.  Su una
          frazione (stimata $\sim$30\%) di trend duraturi, il
          pyramiding contribuisce ulteriori 2-3 unita' di profitto
          rispetto alla baseline.
\end{itemize}

In condizioni nominali ci attendiamo dunque un PnL medio
nell'ordine di $+100\$$ -- $+250\$$ per giornata di trading, con
deviazione standard maggiore della baseline (a causa
dell'amplificazione del pyramiding) ma ratio Sharpe simile o
superiore.  Le verifiche empiriche sul \texttt{tournament.ipynb}
sono in corso al momento della consegna del report.

\paragraph{Limiti e potenziali ottimizzazioni.}
\begin{itemize}
    \item \emph{Calibrazione statica}: i parametri sono fissati a
          valori derivati dall'analisi della volatilita' intraday
          tipica di RMSC04; una variante ATR-adattiva potrebbe
          migliorare la robustezza in regimi di volatilita' anomala.
    \item \emph{Assenza di filtro macro}: non utilizziamo segnali
          a piu' alto timeframe (ad es.\ momentum giornaliero) per
          confermare il trend intraday; un \emph{trend filter}
          giornaliero potrebbe ridurre i falsi segnali nei
          mercati laterali estesi.
    \item \emph{Sensibilita' al seed}: i sistemi trend-following
          sono notoriamente sensibili al \emph{path} dei prezzi;
          un'analisi su un campione piu' ampio di seed (50+) sarebbe
          desiderabile per stimare con maggiore precisione il
          \emph{worst-case drawdown}.
\end{itemize}

% =============================================================
\section{Conclusioni}
\label{sec:conclusions}
% =============================================================

Abbiamo presentato un \emph{trading agent} a strategia
trend-following per il simulatore ABIDES nella configurazione
RMSC04.  Le scelte di design sono state guidate da tre principi:
\textbf{semplicita'} (un singolo file Python di circa 400 righe,
nessuna dipendenza esterna oltre \texttt{numpy} e \texttt{abides\_*}),
\textbf{robustezza} (un \emph{triple stop system} che protegge
gerarchicamente per trade, per sessione e per giornata) e
\textbf{rispondenza alla microstruttura del mercato simulato} (la
mean reversion del fundamental OU di RMSC04 genera microtrend
sfruttabili da una doppia EMA con timeframe brevi).

Lavori futuri comprendono: (i)~l'introduzione di una versione
ATR-adattiva degli stop, (ii)~l'aggiunta di un filtro macro
giornaliero, (iii)~lo sviluppo di una variante con \emph{reinforcement
learning} che apprenda online i parametri ottimali su un orizzonte
multi-giorno.  In una prospettiva piu' ampia, la metodologia
modulare presentata (signal + targeting + pyramiding + triple stop)
e' direttamente trasferibile ad altre famiglie di strategie nello
stesso framework ABIDES.

% =============================================================
\section*{Ringraziamenti}
% =============================================================

Si ringraziano i docenti del corso Algorithmic Trading 2026 della
PoliMi Graduate School of Management per la disponibilita' del
framework ABIDES, dei notebook \texttt{agent\_monitor.ipynb} e
\texttt{tournament.ipynb} e degli agenti esempio
\texttt{BasicAgent}, \texttt{MeanReversionAgent},
\texttt{TrendFollowerAgent} e \texttt{LiquidityProviderAgent}
utilizzati come baseline e codebase di partenza.

% =============================================================
\begin{thebibliography}{9}
% =============================================================

\bibitem{byrd2020abides}
Byrd, D., Hybinette, M., Balch, T.\ H.
\emph{ABIDES: Towards High-Fidelity Market Simulation for AI
Research}.
Proceedings of the ACM International Conference on AI in Finance
(ICAIF), 2020.

\bibitem{murphy1999technical}
Murphy, J.\ J.
\emph{Technical Analysis of the Financial Markets: A Comprehensive
Guide to Trading Methods and Applications}.
New York Institute of Finance, 1999.

\bibitem{wilder1978new}
Wilder, J.\ W.\ Jr.
\emph{New Concepts in Technical Trading Systems}.
Trend Research, 1978.

\bibitem{chande1994new}
Chande, T., Kroll, S.
\emph{The New Technical Trader: Boost Your Profit by Plugging into
the Latest Indicators}.
Wiley, 1994.

\bibitem{covel2007turtle}
Covel, M.\ W.
\emph{The Complete TurtleTrader: The Legend, the Lessons, the
Results}.
HarperBusiness, 2007.

\bibitem{tharp1998trade}
Tharp, V.\ K.
\emph{Trade Your Way to Financial Freedom}.
McGraw-Hill, 1998.

\bibitem{vince1992mathematics}
Vince, R.
\emph{The Mathematics of Money Management: Risk Analysis Techniques
for Traders}.
Wiley, 1992.

\bibitem{bertsimas1998optimal}
Bertsimas, D., Lo, A.\ W.
\emph{Optimal Control of Execution Costs}.
Journal of Financial Markets, vol.\ 1, no.\ 1, pp.\ 1--50, 1998.

\bibitem{avellaneda2010stat}
Avellaneda, M., Lee, J.\ H.
\emph{Statistical Arbitrage in the U.S.\ Equities Market}.
Quantitative Finance, vol.\ 10, no.\ 7, pp.\ 761--782, 2010.

\end{thebibliography}

\end{document}
